% !TEX root = ../main.tex
\chapter{Discusión}

En el presente capítulo se interpretan conjuntamente los resultados de las
cuatro líneas experimentales. A partir de ellos se responde la pregunta de
investigación, se establece su relación con los trabajos previos y se precisan
las principales limitaciones del estudio.

\section{Respuesta a la pregunta de investigación}

La pregunta de investigación planteó si imponer equivarianza a rotaciones
continuas durante el aprendizaje no supervisado producía representaciones más
útiles que las de un VAE convolucional comparable. Los resultados no permiten
responder afirmativamente de manera general. En reconstrucción, las cuatro
estimaciones puntuales favorecieron ligeramente al VAE base y el intervalo del
criterio primario incluyó cero. En diagnóstico de WSI, el
VAE-\(\mathrm{SO}(2)\) obtuvo mejores valores puntuales de macro-F1 y exactitud,
pero el
intervalo por WSI de la diferencia primaria también atravesó cero. Finalmente,
en la segmentación dispersa a resolución de parche, el área bajo la curva de
eficiencia de etiquetas fue mayor para la rama equivariante, aunque su
intervalo simultáneo tampoco excluyó cero.

Por tanto, la hipótesis inicial de una mejora estadísticamente concluyente en
las tareas posteriores no quedó validada. Esta conclusión no equivale a probar
que ambos modelos sean equivalentes ni a demostrar ausencia de beneficio. La
prueba sellada incluyó 23 WSI, algunas clases estuvieron representadas por solo
dos o cinco láminas y cada rama supervisada se entrenó una sola vez. Los
intervalos reflejan el remuestreo de WSI observadas, pero no la variación entre
inicializaciones ni entre conjuntos de datos. Con ese alcance, la evidencia
impide sostener una afirmación fuerte de superioridad, pero no descarta toda
diferencia posible entre las representaciones.

Sí se verificó una propiedad más específica. Al aplicar al latente la rotación
prescrita y decodificar, el error fue menor para el VAE-\(\mathrm{SO}(2)\) en
los 25 parches del conjunto fijo. En cuartos de vuelta, la razón de acción
descendió de 0,709453 a \(2,18\times10^{-6}\), y las reconstrucciones del
modelo base mostraron artefactos que no aparecieron en la rama equivariante.
Para rotaciones interpoladas cada cinco grados, las medianas de acción y
canonización también disminuyeron, aunque no alcanzaron el umbral absoluto
predefinido. En consecuencia, los experimentos sí muestran que la restricción
arquitectónica se conserva como una regularidad verificable en la
representación latente y en la decodificación; no muestran todavía una
equivarianza continua perfecta de extremo a extremo.

\section{Comparación entre representaciones}

La reconstrucción revela un primer compromiso. Restringir los núcleos redujo
el número de parámetros libres, pero amplió sus soportes y aumentó el costo de
convolución estimado. El VAE base obtuvo un MAE medio ligeramente menor, junto
con mejores MSE, PSNR y SSIM descriptivos. Esto indica que la restricción
geométrica no produjo una mejora gratuita de fidelidad. No obstante, la
diferencia primaria fue pequeña y el experimento no fue una prueba de
equivalencia o no inferioridad; por ello tampoco puede concluirse que la
calidad de reconstrucción de una rama sea intercambiable con la de la otra.

Las tareas supervisadas sugieren una interacción con la cantidad de etiquetas.
En diagnóstico, la diferencia puntual de macro-F1 fue de 0,095 a favor del
VAE-\(\mathrm{SO}(2)\), pero con un intervalo amplio. En tejido, la rama
equivariante obtuvo mayor macro-F1 con 250, 500, \(1\,000\) y \(2\,500\)
etiquetas por clase; solo en el presupuesto de 500 el intervalo simultáneo
excluyó cero. Con todas las \(5\,671\) etiquetas disponibles, la estimación
puntual se invirtió y favoreció al VAE base. La secuencia es compatible con la
posibilidad de que el sesgo geométrico ayude en un régimen de supervisión muy
escasa y pierda importancia cuando el predictor dispone de más ejemplos. Sin
embargo, la falta de monotonía, el tamaño de los intervalos y la única
inicialización impiden presentarla como una conclusión general sobre la
eficiencia muestral.

En este sentido, la lectura anterior puede relacionarse con la \emph{Bitter
Lesson} de Sutton~\cite{sutton_bitter_lesson_2019}. Su argumento histórico es
que, a largo plazo, los métodos generales capaces de aprovechar más cómputo y
datos suelen superar soluciones que incorporan conocimiento humano específico.
La equivarianza empleada en esta tesis constituye precisamente un sesgo
incorporado, pues restringe de antemano cuáles transformaciones puede aprender
cada capa. La desaparición de la ventaja puntual al usar todas las etiquetas es
compatible con esa tensión, pero el presente experimento no constituye una
prueba de la tesis de Sutton. No se variaron de manera sistemática la escala
del modelo, el cómputo y el volumen de preentrenamiento y, en patología, la
disponibilidad de etiquetas expertas no crece necesariamente al ritmo del
cómputo. En un régimen fijo y pequeño, un sesgo adecuado todavía puede ser
útil; determinar cuándo deja de serlo exige curvas de escalamiento diseñadas
para responder esa pregunta.

\section{Relación con trabajos previos}

Los resultados geométricos deben interpretarse distinguiendo invariancia de
equivarianza. Una representación invariante intenta conservar el mismo vector
cuando rota la entrada. Una representación equivariante permite que el vector
cambie, pero exige que lo haga mediante una acción conocida. Elphick et al.
evaluaron la invariancia a rotaciones de doce modelos fundacionales de
patología mediante vecinos mutuos y distancia coseno, y encontraron mayor
invariancia en los modelos que habían usado aumento por rotación durante el
preentrenamiento~\cite{elphick_rotation_2024}. Ese estudio y esta tesis no
miden la misma propiedad: allí se compara la alineación entre representaciones
globales; aquí se transforma explícitamente un mapa latente y se comprueba su
decodificación. Ambos resultados coinciden, sin embargo, en que la respuesta a
la orientación no debe suponerse a partir del desempeño final, sino medirse de
forma directa.

La disminución del error de acción y la estabilidad visual del
VAE-\(\mathrm{SO}(2)\) aportan evidencia complementaria a
EQ-VAE~\cite{kouzelis_eqvae_2025}. En ambos casos se estudia si una transformación
semánticamente conservativa adquiere una estructura más regular en el espacio
latente. EQ-VAE induce esa propiedad mediante regularización sobre un
autoencodificador preentrenado; en esta tesis se incorporó mediante el espacio
de núcleos direccionables desde el inicio del entrenamiento. La comparación es
conceptual, pues no se reprodujo su pérdida ni su tarea generativa.

La geometría local ofrece una señal más débil. Con pasos de un grado, la medida
corregida disminuyó 7,38\% y favoreció al modelo equivariante en los 25 parches,
pero no alcanzó el margen predefinido de 10\% ni conservó la dirección con
pasos de cinco grados. Asimismo, la medida espacial y el sondeo RGB no mostraron
una organización uniformemente más simple. Con el alcance de estas pruebas no
es posible afirmar que todo el espacio latente sea suave o esté globalmente
factorizado; lo que sí se observa es que una acción rotacional concreta puede
ejecutarse con mucho menor error sobre la representación equivariante.

\section{Limitaciones y amenazas a la validez}

La principal limitación estadística es el número de unidades biológicas
independientes. Aunque la reconstrucción incluyó \(67\,138\) parches, estos
provenían de solo 23 WSI de prueba y no deben contarse como decenas de miles de
réplicas independientes. El bootstrap por conglomerados respetó esa estructura,
pero no compensa la variación biológica que el conjunto no contiene. La baja
presencia de LGSC, MC y necrosis amplió la incertidumbre y limita las
conclusiones por clase.

La segunda limitación es la variación experimental. Se comparó una
inicialización por arquitectura y, por tanto, no fue posible separar el efecto
del sesgo inductivo de la variación entre entrenamientos. Además de la
restricción de simetría, las ramas difirieron en parametrización, número de
parámetros y costo computacional. Las posiciones de las no linealidades y los
clasificadores posteriores se igualaron, pero la comparación no equivale a un
experimento controlado en el que solo cambie el sesgo de simetría.

El análisis rotacional también tiene límites propios. Los cuartos de vuelta
son reindexaciones exactas de la cuadrícula y verifican la compatibilidad con
\(C_4\), mientras los demás ángulos requieren interpolación de imágenes y
campos. Las 360 vistas de un parche son mediciones repetidas, no nuevas muestras
biológicas. El conjunto fijo de 25 parches permitió comparaciones pareadas y
visuales controladas, pero no basta para generalizar la geometría observada a
todo UBC-OCEAN ni a otros órganos, tinciones o centros clínicos.

Por último, la segmentación se evaluó solo en regiones anotadas de alta pureza.
Las máscaras suplementarias no permiten valorar exhaustivamente la WSI y sus
zonas negras no constituyen negativos. Tampoco se conservaron mapas de atención
de la inferencia sellada. Quedan pendientes mediciones más amplias de
factorización, topología y acción global; por ello las conclusiones se
restringen a lo ya observado.
